% !TEX root = ../../../main.tex
% 所属：第一部分：单变量分析学 / 第10章：一些定积分的问题
% 节：10.3 积分不等式~~ 一
% 正文按原文件顺序保存；由 ../chapter10.tex 调用。
\section{积分不等式~~ 一}

本节主要整理定积分不等式的若干证明方法：构造辅助函数、利用被积函数的单调性与对称性、逐次积分得到不等式，以及通过分部积分选择合适的辅助因子。

\subsection*{用辅助函数方法证明积分不等式}

\begin{problem}
	设 $f(x)$ 在 $[0,1]$ 上可微，且当 $x\in(0,1)$ 时，
	\[
	0<f'(x)<1,\qquad f(0)=0.
	\]
	证明：
	\[
	\left(\int_0^1 f(x)\dd x\right)^2>\int_0^1 f^3(x)\dd x.
	\]
\end{problem}

\begin{solution}
	由 $f(0)=0$ 与 $f'(x)>0$ 知，当 $x\in(0,1)$ 时有 $f(x)>0$。
	
	\textbf{证法一.} 令
	\[
	F(x)=\left(\int_0^x f(t)\dd t\right)^2-\int_0^x f^3(t)\dd t.
	\]
	则 $F(0)=0$，且
	\[
	F'(x)
	=2f(x)\int_0^x f(t)\dd t-f^3(x)
	=f(x)\left(2\int_0^x f(t)\dd t-f^2(x)\right).
	\]
	再令
	\[
	G(x)=2\int_0^x f(t)\dd t-f^2(x).
	\]
	则 $G(0)=0$，且
	\[
	G'(x)=2f(x)-2f(x)f'(x)=2f(x)(1-f'(x))>0.
	\]
	故 $G(x)>0$，从而 $F'(x)>0$，于是
	\[
	F(1)>F(0)=0.
	\]
	即
	\[
	\left(\int_0^1 f(x)\dd x\right)^2>\int_0^1 f^3(x)\dd x.
	\]
	
	\textbf{证法二.} 令
	\[
	F(x)=\left(\int_0^x f(t)\dd t\right)^2,\qquad
	G(x)=\int_0^x f^3(t)\dd t.
	\]
	由 Cauchy 中值定理，存在 $\xi\in(0,1)$，使得
	\[
	\frac{\left(\int_0^1 f(x)\dd x\right)^2}{\int_0^1 f^3(x)\dd x}
	=\frac{F(1)-F(0)}{G(1)-G(0)}
	=\frac{F'(\xi)}{G'(\xi)}
	=\frac{2\int_0^\xi f(t)\dd t}{f^2(\xi)}.
	\]
	再对
	\[
	H(x)=2\int_0^x f(t)\dd t,\qquad K(x)=f^2(x)
	\]
	在 $[0,\xi]$ 上用 Cauchy 中值定理，存在 $\eta\in(0,\xi)$，使得
	\[
	\frac{2\int_0^\xi f(t)\dd t}{f^2(\xi)}
	=\frac{H(\xi)-H(0)}{K(\xi)-K(0)}
	=\frac{2f(\eta)}{2f(\eta)f'(\eta)}
	=\frac1{f'(\eta)}>1.
	\]
	故原不等式成立。
\end{solution}

\begin{problem}
	设 $f(x)\in C[0,1]$，且满足 $0\leq f(x)\leq x$。证明：
	\[
	\int_0^1 x^2 f(x)\dd x\geq \left(\int_0^1 f(x)\dd x\right)^2,
	\]
	并求使上式成为等式的所有连续函数 $f(x)$。
\end{problem}

\begin{solution}
	令
	\[
	F(x)=\int_0^x t^2 f(t)\dd t-\left(\int_0^x f(t)\dd t\right)^2.
	\]
	则 $F(0)=0$，且
	\[
	F'(x)=x^2f(x)-2f(x)\int_0^x f(t)\dd t
	=f(x)\left(x^2-2\int_0^x f(t)\dd t\right).
	\]
	令
	\[
	G(x)=x^2-2\int_0^x f(t)\dd t.
	\]
	由于
	\[
	G(0)=0,\qquad
	G'(x)=2x-2f(x)\geq 0,
	\]
	故 $G(x)\geq0$。又 $f(x)\geq0$，于是 $F'(x)\geq0$，所以
	\[
	F(1)\geq F(0)=0,
	\]
	即
	\[
	\int_0^1 x^2 f(x)\dd x\geq \left(\int_0^1 f(x)\dd x\right)^2.
	\]
	
	下面讨论等号情形。若 $f\equiv0$，显然等号成立；若 $f\not\equiv0$，则存在 $c\in(0,1]$ 使 $f(c)>0$。等号成立时 $F(1)=0$，又 $F'\geq0$，故 $F'\equiv0$。于是 $f(c)G(c)=0$，从而 $G(c)=0$。但
	\[
	G(c)=2\int_0^c(t-f(t))\dd t,
	\]
	且 $t-f(t)\geq0$，故 $f(t)=t$ 在 $[0,c]$ 上成立。由连续性再向右延拓可得 $f(x)=x$ 在 $[0,1]$ 上成立。
	
	因此等号成立当且仅当
	\[
	f(x)\equiv0\qquad\text{或}\qquad f(x)=x.
	\]
\end{solution}

\begin{problem}
	设 $f(x),g(x)$ 为 $[a,b]$ 上单调递增的连续函数，证明：
	\[
	(b-a)\int_a^b f(x)g(x)\dd x
	\geq
	\int_a^b f(x)\dd x\int_a^b g(x)\dd x.
	\]
\end{problem}

\begin{solution}
	\[
	F(x)=(x-a)\int_a^x f(t)g(t)\dd t
	-
	\int_a^x f(t)\dd t\int_a^x g(t)\dd t.
	\]
	\[F'(x) = \int_{a}^{x}(f(t)-f(x))(g(t)-g(x))\dd t>0\]
	又$F(0)=0$，得证
\end{solution}

\begin{problem}
	设函数 $f(x),g(x)$ 在 $[a,b]$ 上连续，且 $f(x)$ 单调增加，$0\leq g(x)\leq1$。证明：
	\[
	\int_a^{a+\int_a^b g(t)\dd t} f(x)\dd x
	\leq
	\int_a^b f(x)g(x)\dd x.
	\]
\end{problem}

\begin{solution}
	\[
	F(x)=\int_a^{a+\int_a^x g(t)\dd t} f(t)\dd t
	-
	\int_a^x f(t)g(t)\dd t.
	\]
	\[F'(x)=f(a+\int_a^b g(t)\dd t)-f(x)g(x)\]
	由$0\leq g(x)\leq1$，$\displaystyle\int_{a}^{x}g(x)\dd x    \leqslant x-a$,且 $f(x)$ 单调增加，故\[f(a+\int_a^b g(t)\dd t)-f(x)g(x)\leqslant 0\]
	从而$F'(x)\leqslant 0$\[F(x)\leqslant 0\]
\end{solution}

\subsection*{利用被积函数的性质}

\begin{problem}
	证明：
	\[
	\int_0^1\frac{\cos x}{\sqrt{1-x^2}}\dd x>
	\int_0^1\frac{\sin x}{\sqrt{1-x^2}}\dd x.
	\]
\end{problem}

\begin{solution}
	对左端作代换 $x=\sin t$，得
	\[
	\int_0^1\frac{\cos x}{\sqrt{1-x^2}}\dd x
	=\int_0^{\frac\pi2}\cos(\sin t)\dd t.
	\]
	对右端作代换 $x=\cos t$，得
	\[
	\int_0^1\frac{\sin x}{\sqrt{1-x^2}}\dd x
	=\int_0^{\frac\pi2}\sin(\cos t)\dd t.
	\]
	而对 $t\in[0,\pi/2]$，有
	\[
	\frac\pi2-\sin t>\cos t,
	\]
	且二者均在 $[0,\pi/2]$ 内。由 $\sin x$ 在 $[0,\pi/2]$ 上单调递增，得
	\[
	\cos(\sin t)>\cos t>\sin(\cos t).
	\]
	两端积分，即得所证。
\end{solution}

\begin{problem}
	函数 $f(x)$ 在闭区间 $[0,1]$ 上有连续的一阶导数，证明：
	\[
	\int_0^1 |f(x)|\dd x\leq
	\max\left\{\int_0^1 |f'(x)|\dd x,\ \left|\int_0^1 f(x)\dd x\right|\right\}.
	\]
\end{problem}

\begin{solution}
	若
	\[
	\left|\int_0^1 f(x)\dd x\right|=\int_0^1 |f(x)|\dd x,
	\]
	则结论显然成立。
	
	若
	\[
	\left|\int_0^1 f(x)\dd x\right|<\int_0^1 |f(x)|\dd x,
	\]
	则 $f$ 在 $[0,1]$ 上变号。由连续性，存在 $x_0\in(0,1)$，使得 $f(x_0)=0$。于是对任意 $x\in[0,1]$，
	\[
	|f(x)|=|f(x)-f(x_0)|
	=\left|\int_{x_0}^x f'(t)\dd t\right|
	\leq\int_0^1 |f'(t)|\dd t.
	\]
	积分即得
	\[
	\int_0^1 |f(x)|\dd x\leq \int_0^1 |f'(x)|\dd x.
	\]
	故原不等式成立。
\end{solution}

\begin{problem}
	函数 $f(x)$ 在 $[0,1]$ 上单调不增，证明：对于任意 $\alpha\in(0,1)$，
	\[
	\int_0^\alpha f(x)\dd x\geq \alpha\int_0^1 f(x)\dd x.
	\]
\end{problem}

\begin{solution}
	原不等式等价于
	\[
	(1-\alpha)\int_0^\alpha f(x)\dd x
	\geq
	\alpha\int_\alpha^1 f(x)\dd x,
	\]
	也即
	\[
	\frac1\alpha\int_0^\alpha f(x)\dd x
	\geq
	\frac1{1-\alpha}\int_\alpha^1 f(x)\dd x.
	\]
	由积分中值定理，存在 $\xi\in(0,\alpha)$ 与 $\eta\in(\alpha,1)$，使得
	\[
	\frac1\alpha\int_0^\alpha f(x)\dd x=f(\xi),\qquad
	\frac1{1-\alpha}\int_\alpha^1 f(x)\dd x=f(\eta).
	\]
	由于 $f$ 单调不增，且 $\xi<\eta$，故 $f(\xi)\geq f(\eta)$。结论成立。
\end{solution}

\begin{problem}
	设函数 $f(x)\geq0$ 在 $[0,1]$ 上连续递减，试证：若 $0<\alpha<\beta<1$，则
	\[
	\int_0^\alpha f(x)\dd x
	\geq
	\frac{\alpha}{\beta-\alpha}\int_\alpha^\beta f(x)\dd x
	>
	\frac{\alpha}{\beta}\int_0^\beta f(x)\dd x.
	\]
\end{problem}

\begin{solution}
	由积分中值定理，存在 $\xi\in(0,\alpha)$ 与 $\eta\in(\alpha,\beta)$，使得
	\[
	\int_0^\alpha f(x)\dd x=\alpha f(\xi),\qquad
	\frac{\alpha}{\beta-\alpha}\int_\alpha^\beta f(x)\dd x=\alpha f(\eta).
	\]
	由于 $f$ 递减，且 $\xi<\eta$，故 $f(\xi)\geq f(\eta)$，从而
	\[
	\int_0^\alpha f(x)\dd x
	\geq
	\frac{\alpha}{\beta-\alpha}\int_\alpha^\beta f(x)\dd x.
	\]

	所以
\[
\int_0^\alpha f(x)\dd x
\geq
\frac{\alpha}{\beta-\alpha}\int_\alpha^\beta f(x)\dd x
>
\frac{\alpha}{\beta}\int_0^\beta f(x)\dd x.
\]
\end{solution}

\begin{problem}
	设 $f(x)$ 是 $[a,b]$ 上连续递增函数，证明：
	\[
	\int_a^b x f(x)\dd x
	\geq
	\frac{a+b}{2}\int_a^b f(x)\dd x.
	\]
\end{problem}

\begin{solution}
	记
	\[
	c=\frac{a+b}{2},\qquad h=\frac{b-a}{2}.
	\]
	只需证明
	\[
	\int_a^b (x-c)f(x)\dd x\geq0.
	\]
	作变换 $x=c-t$ 与 $x=c+t$，得
	\[
	\begin{aligned}
		\int_a^b (x-c)f(x)\dd x
		&=-\int_0^h t f(c-t)\dd t+\int_0^h t f(c+t)\dd t\\
		&=\int_0^h t\bigl(f(c+t)-f(c-t)\bigr)\dd t.
	\end{aligned}
	\]
	由于 $f$ 递增，故 $f(c+t)\geq f(c-t)$，于是上式非负，结论成立。
\end{solution}

\begin{problem}
	设 $a,b>0,\ f(x)\geq0$，且 $f(x)$ 在 $[a,b]$ 上可积，
	\[
	\int_{a}^b x f(x)\dd x=0.
	\]
	证明：
	\[
	\int_{a}^b x^2 f(x)\dd x\leq
	ab\int_{a}^b f(x)\dd x.
	\]
\end{problem}

\begin{solution}
	当 $x\in[a,b]$ 时，
	\[
	(x+a)(b-x)f(x)\geq0.
	\]
	两端积分得
	\[
	\int_{a}^b\bigl(-x^2+(b-a)x+ab\bigr)f(x)\dd x\geq0.
	\]
	由
	\[
	\int_{a}^b x f(x)\dd x=0
	\]
	可得
	\[
	-\int_{a}^b x^2 f(x)\dd x
	+ab\int_{a}^b f(x)\dd x\geq0,
	\]
	即所证。
\end{solution}

\begin{problem}
	设 $f(x)$ 与 $g(x)$ 为 $[a,b]$ 上同序的，即任意 $x_1,x_2\in[a,b]$ 有
	\[
	(f(x_1)-f(x_2))(g(x_1)-g(x_2))\geq0.
	\]
	证明：
	\[
	\int_a^b f(x)\dd x\int_a^b g(x)\dd x
	\leq
	(b-a)\int_a^b f(x)g(x)\dd x.
	\]
\end{problem}

\begin{solution}
	将
	\[
	(f(x)-f(y))(g(x)-g(y))\geq0
	\]
	在 $[a,b]\times[a,b]$ 上积分，得
	\[
	0\leq
	\int_a^b\int_a^b(f(x)-f(y))(g(x)-g(y))\dd x\dd y.
	\]
	展开得
	\[
	0\leq
	2(b-a)\int_a^b f(x)g(x)\dd x
	-2\int_a^b f(x)\dd x\int_a^b g(x)\dd x.
	\]
	整理即得结论。
\end{solution}

\begin{problem}
	设 $f'(x)$ 在 $[a,b]$ 上连续。证明：
	\[
	\max_{a\leq x\leq b}|f(x)|
	\leq
	\left|\frac1{b-a}\int_a^b f(x)\dd x\right|
	+\int_a^b |f'(x)|\dd x.
	\]
\end{problem}

\begin{solution}
	由积分中值定理，存在 $\xi\in[a,b]$，使得
	\[
	f(\xi)=\frac1{b-a}\int_a^b f(x)\dd x.
	\]
	任取 $x\in[a,b]$，则
	\[
	\begin{aligned}
		|f(x)|
		&=|f(x)-f(\xi)+f(\xi)|\\
		&\leq \left|\int_\xi^x f'(t)\dd t\right|+|f(\xi)|\\
		&\leq \int_a^b |f'(t)|\dd t
		+\left|\frac1{b-a}\int_a^b f(t)\dd t\right|.
	\end{aligned}
	\]
	对 $x$ 取最大值，即得所证。
\end{solution}

\begin{problem}
	设 $f(x)$ 在 $[0,1]$ 上二阶连续可导，
	\[
	f(0)=f(1)=0,\qquad f(x)\neq0\quad (\forall x\in(0,1)),
	\]
	且
	\[
	\int_0^1\left|\frac{f''(x)}{f(x)}\right|\dd x
	\]
	存在有限。证明：
	\[
	\int_0^1\left|\frac{f''(x)}{f(x)}\right|\dd x\geq4.
	\]
\end{problem}

\begin{solution}
	由条件知 $f$ 在 $(0,1)$ 内不变号。取 $x_0\in(0,1)$，使得
	\[
	|f(x_0)|=\max_{0\leq x\leq1}|f(x)|.
	\]
	则 $f'(x_0)=0$，且 $|f(x_0)|>0$。
	
	由 Lagrange 中值定理，存在 $\xi\in(0,x_0)$ 与 $\eta\in(x_0,1)$，使得
	\[
	f'(\xi)=\frac{f(x_0)-f(0)}{x_0}
	=\frac{f(x_0)}{x_0},
	\]
	\[
	f'(\eta)=\frac{f(1)-f(x_0)}{1-x_0}
	=-\frac{f(x_0)}{1-x_0}.
	\]
	于是
	\[
	\begin{aligned}
		\int_0^1 |f''(x)|\dd x
		&\geq \left|\int_\xi^\eta f''(x)\dd x\right|\\
		&=|f'(\eta)-f'(\xi)|\\
		&=|f(x_0)|\left(\frac1{x_0}+\frac1{1-x_0}\right)\\
		&\geq4|f(x_0)|.
	\end{aligned}
	\]
	又 $|f(x)|\leq|f(x_0)|$，故
	\[
	\int_0^1\left|\frac{f''(x)}{f(x)}\right|\dd x
	=
	\int_0^1\frac{|f''(x)|}{|f(x)|}\dd x
	\geq
	\frac1{|f(x_0)|}\int_0^1 |f''(x)|\dd x
	\geq4.
	\]
\end{solution}

\begin{problem}
	设函数 $f(x):[0,1]\to\mathbb{R}$ 连续可微，并且
	\[
	\int_0^1 f(x)\dd x=0.
	\]
	证明：对于任意 $a\in(0,1)$，有
	\[
	\left|\int_0^a f(x)\dd x\right|
	\leq
	\frac18\max_{0\leq x\leq1}|f'(x)|.
	\]
\end{problem}

\begin{solution}
	令
	\[
	g(x)=\int_0^x f(t)\dd t.
	\]
	则
	\[
	g(0)=g(1)=0,\qquad g'(x)=f(x),\qquad g''(x)=f'(x).
	\]
	记
	\[
	M=\max_{0\leq x\leq1}|g''(x)|=\max_{0\leq x\leq1}|f'(x)|.
	\]
	对 $g$ 分别在点 $a$ 处向 $0$ 与 $1$ 作 Taylor 公式，存在 $\xi_1\in(0,a)$ 与 $\xi_2\in(a,1)$，使得
	\[
	0=g(0)=g(a)-ag'(a)+\frac{g''(\xi_1)}2a^2,
	\]
	\[
	0=g(1)=g(a)+(1-a)g'(a)+\frac{g''(\xi_2)}2(1-a)^2.
	\]
	第一式乘以 $1-a$，第二式乘以 $a$，相加得
	\[
	g(a)=-\frac{a(1-a)}2\left(ag''(\xi_1)+(1-a)g''(\xi_2)\right).
	\]
	故
	\[
	|g(a)|\leq \frac{a(1-a)}2M\leq\frac18M.
	\]
	即
	\[
	\left|\int_0^a f(x)\dd x\right|
	\leq
	\frac18\max_{0\leq x\leq1}|f'(x)|.
	\]
\end{solution}

\begin{problem}
	设函数 $f(x)$ 在 $[a,b]$ 上连续可导。证明：
	\[
	\int_a^b |f(x)|\dd x
	\leq
	\max\left\{
	(b-a)\int_a^b |f'(x)|\dd x,\ 
	\left|\int_a^b f(x)\dd x\right|
	\right\}.
	\]
\end{problem}

\begin{solution}
	若
	\[
	\int_a^b |f(x)|\dd x=\left|\int_a^b f(x)\dd x\right|,
	\]
	则结论显然成立。
	
	否则 $f$ 在 $[a,b]$ 上变号。由连续性，存在 $x_0\in[a,b]$，使得 $f(x_0)=0$。于是对任意 $x\in[a,b]$，
	\[
	|f(x)|=|f(x)-f(x_0)|
	=\left|\int_{x_0}^x f'(t)\dd t\right|
	\leq
	\int_a^b |f'(t)|\dd t.
	\]
	两端积分得
	\[
	\int_a^b |f(x)|\dd x
	\leq
	(b-a)\int_a^b |f'(x)|\dd x.
	\]
	故原不等式成立。
\end{solution}

\subsection*{对三角函数不等式的两例练习}

\begin{problem}
	证明：当 $x>0$ 时，
	\[
	x-\frac{x^3}{6}<\sin x<x-\frac{x^3}{6}+\frac{x^5}{120}.
	\]
\end{problem}

\begin{solution}
	由于 $\cos x\leq1$，对 $[0,x]$ 积分得
	\[
	\sin x=\int_0^x\cos t\dd t\leq\int_0^x\dd t=x.
	\]
	再积分一次，
	\[
	1-\cos x=\int_0^x\sin t\dd t\leq\int_0^x t\dd t=\frac{x^2}{2}.
	\]
	于是
	\[
	x-\sin x=\int_0^x(1-\cos t)\dd t
	\leq\int_0^x\frac{t^2}{2}\dd t=\frac{x^3}{6},
	\]
	故
	\[
	x-\frac{x^3}{6}\leq\sin x.
	\]
	由于 $x>0$，上述不等式实际上为严格不等式。
	
	另一方面，由
	\[
	x-\sin x\leq\frac{x^3}{6}
	\]
	再积分，得
	\[
	\frac{x^2}{2}+\cos x-1
	=\int_0^x(t-\sin t)\dd t
	\leq\int_0^x\frac{t^3}{6}\dd t=\frac{x^4}{24}.
	\]
	继续积分，得
	\[
	\frac{x^3}{6}-x+\sin x
	=\int_0^x\left(\frac{t^2}{2}-1+\cos t\right)\dd t
	\leq\int_0^x\frac{t^4}{24}\dd t
	=\frac{x^5}{120}.
	\]
	故
	\[
	\sin x\leq x-\frac{x^3}{6}+\frac{x^5}{120}.
	\]
	同样因为 $x>0$，上式为严格不等式。
\end{solution}

\subsection*{用待定系数法证明}

\begin{problem}
	设 $f(x)$ 在 $[0,1]$ 上有一阶连续导数，且 $f(0)=f(1)=0$，
	\[
	M=\max_{x\in[0,1]}|f'(x)|.
	\]
	证明：
	\[
	\left|\int_0^1 f(x)\dd x\right|\leq \frac{M}{4}.
	\]
\end{problem}

\begin{solution}
	对任意 $a\in[0,1]$，由分部积分得
	\[
	\int_0^1 f(x)\dd x
	=\int_0^1 f(x)\dd(x-a)
	=[f(x)(x-a)]_0^1-\int_0^1 f'(x)(x-a)\dd x
	=-\int_0^1 f'(x)(x-a)\dd x.
	\]
	于是
	\[
	\left|\int_0^1 f(x)\dd x\right|
	\leq
	M\int_0^1|x-a|\dd x
	=M\left(\int_0^a(a-x)\dd x+\int_a^1(x-a)\dd x\right).
	\]
	所以
	\[
	\left|\int_0^1 f(x)\dd x\right|
	\leq
	M\left(a^2-a+\frac12\right).
	\]
	取 $a=\dfrac12$，即得
	\[
	\left|\int_0^1 f(x)\dd x\right|\leq \frac{M}{4}.
	\]
\end{solution}

\begin{problem}
	设函数 $f(x)$ 在闭区间 $[0,1]$ 上连续，且
	\[
	\int_0^1 f(x)\dd x=0,\qquad
	\int_0^1 x f(x)\dd x=1.
	\]
	求证：
	\[
	(1)\quad \exists x_0\in(0,1),\ |f(x_0)|>4;
	\]
	\[
	(2)\quad \exists x_1\in[0,1],\ |f(x_1)|=4.
	\]
\end{problem}

\begin{solution}
	由 $\int_0^1 f(x)\dd x=0$，得
	\[
	1=\int_0^1 x f(x)\dd x
	=\int_0^1\left(x-\frac12\right)f(x)\dd x.
	\]
	若 $|f(x)|\leq4$ 在 $[0,1]$ 上恒成立，则
	\[
	1
	=\left|\int_0^1\left(x-\frac12\right)f(x)\dd x\right|
	\leq
	4\int_0^1\left|x-\frac12\right|\dd x
	=1.
	\]
	若等号成立，则三角不等式与放缩中都必须取等号，从而 $f(x)$ 的符号须与 $x-\frac12$ 的符号一致，且在除 $x=\frac12$ 外几乎处处满足 $|f(x)|=4$。这与 $f$ 的连续性矛盾。因此必有一点 $x_0\in(0,1)$，使得
	\[
	|f(x_0)|>4.
	\]
	又由 $\int_0^1 f(x)\dd x=0$ 以及 $\int_0^1 x f(x)\dd x=1$ 知 $f$ 不恒为 $0$，且 $f$ 在 $[0,1]$ 上必有零点或至少有一点满足 $|f(x)|<4$。由连续性与介值定理，存在 $x_1\in[0,1]$，使得
	\[
	|f(x_1)|=4.
	\]
\end{solution}

\begin{problem}
	设 $f\in C^1[0,1]$，且 $f(0)=f(1)=0$。证明：
	\[
	\left(\int_0^1 f(x)\dd x\right)^2
	\leq
	\frac1{12}\int_0^1 (f'(x))^2\dd x.
	\]
\end{problem}

\begin{solution}
	对任意 $a\in[0,1]$，由分部积分得
	\[
	\int_0^1 f(x)\dd x
	=-\int_0^1 (x-a)f'(x)\dd x.
	\]
	由 Cauchy 不等式(下一节讲)，
	\[
	\left(\int_0^1 f(x)\dd x\right)^2
	\leq
	\int_0^1(x-a)^2\dd x
	\int_0^1(f'(x))^2\dd x.
	\]
	取 $a=\dfrac12$，则
	\[
	\int_0^1\left(x-\frac12\right)^2\dd x=\frac1{12}.
	\]
	故
	\[
	\left(\int_0^1 f(x)\dd x\right)^2
	\leq
	\frac1{12}\int_0^1(f'(x))^2\dd x.
	\]
\end{solution}

\begin{problem}
	设 $f(x)$ 在区间 $[0,1]$ 上连续可导，且 $f(0)=f(1)=0$。求证：
	\[
	\left(\int_0^1 x f(x)\dd x\right)^2
	\leq
	\frac1{45}\int_0^1 (f'(x))^2\dd x,
	\]
	等号当且仅当
	\[
	f(x)=A(x-x^3)
	\]
	时成立，其中 $A$ 是常数。
\end{problem}

\begin{solution}
	对任意常数 $a$，由分部积分得
	\[
	\int_0^1 x f(x)\dd x
	=\frac12\int_0^1 f(x)\dd(x^2-a)
	=-\frac12\int_0^1 (x^2-a)f'(x)\dd x.
	\]
	由 Cauchy 不等式，
	\[
	\left(\int_0^1 x f(x)\dd x\right)^2
	\leq
	\frac14\int_0^1(x^2-a)^2\dd x
	\int_0^1(f'(x))^2\dd x.
	\]
	而
	\[
	\int_0^1(x^2-a)^2\dd x
	=\frac15-\frac{2a}{3}+a^2.
	\]
	右端关于 $a$ 的二次式在 $a=\dfrac13$ 处取最小值
	\[
	\frac15-\frac29+\frac19=\frac4{45}.
	\]
	因此
	\[
	\left(\int_0^1 x f(x)\dd x\right)^2
	\leq
	\frac1{45}\int_0^1(f'(x))^2\dd x.
	\]
	
	等号成立当且仅当 Cauchy 不等式取等号，即存在常数 $\lambda$，使得
	\[
	f'(x)=\lambda\left(x^2-\frac13\right).
	\]
	积分得
	\[
	f(x)=\lambda\left(\frac{x^3}{3}-\frac{x}{3}\right)+C.
	\]
	由 $f(0)=0$ 得 $C=0$，故
	\[
	f(x)=A(x-x^3).
	\]
	反之，此类函数使上面的 Cauchy 不等式取等号，因此等号成立。
\end{solution}
\subsection*{Taylor 公式在积分不等式中的应用}

\begin{problem}
	设 $f(x)$ 在 $[0,a]$ 上二阶可导，$a>0$，且 $f''(x)\geqslant0$，证明：
	\[
	\int_0^a f(x)\dd x\geqslant a f\left(\frac a2\right).
	\]
\end{problem}

\begin{proof}
	对 $x$ 在 $a/2$ 处作 Taylor 展开，由 $f''\geqslant0$ 得
	\[
	f(x)\geqslant f\left(\frac a2\right)+f'\left(\frac a2\right)\left(x-\frac a2\right).
	\]
	两边在 $[0,a]$ 上积分，得
	\[
	\int_0^a f(x)\dd x
	\geqslant a f\left(\frac a2\right)+f'\left(\frac a2\right)\int_0^a\left(x-\frac a2\right)\dd x
	=a f\left(\frac a2\right).
	\]
\end{proof}

\begin{problem}
	设函数 $f(x)$ 在 $[0,1]$ 上二阶可导，且 $f''(x)\leqslant0$，$x\in[0,1]$，证明：
	\[
	\int_0^1 f(x^2)\dd x\leqslant f\left(\frac13\right).
	\]
\end{problem}

\begin{proof}
	对 $f$ 在 $1/3$ 处作 Taylor 展开，由 $f''\leqslant0$ 得
	\[
	f(x^2)\leqslant f\left(\frac13\right)+f'\left(\frac13\right)\left(x^2-\frac13\right).
	\]
	两边在 $[0,1]$ 上积分，注意 $\displaystyle\int_0^1\left(x^2-\frac13\right)\dd x=0$，故
	\[
	\int_0^1 f(x^2)\dd x\leqslant f\left(\frac13\right).
	\]
\end{proof}

\begin{problem}
	设 $f(x)$ 在 $[0,a^n]$ 上二阶可导，$a>0$，且 $f''(x)\geqslant0$，证明：对 $n\geqslant1$，
	\[
	\int_0^a f(x^n)\dd x\geqslant a f\left(\frac{a^n}{n+1}\right).
	\]
\end{problem}

\begin{proof}
	记 $\mu=\dfrac{a^n}{n+1}$。由 $f''\geqslant0$，在 $\mu$ 处作 Taylor 展开得
	\[
	f(x^n)\geqslant f(\mu)+f'(\mu)(x^n-\mu).
	\]
	两边积分，得
	\[
	\int_0^a f(x^n)\dd x
	\geqslant af(\mu)+f'(\mu)\left(\int_0^a x^n\dd x-a\mu\right)
	=af(\mu),
	\]
	即
	\[
	\int_0^a f(x^n)\dd x\geqslant a f\left(\frac{a^n}{n+1}\right).
	\]
\end{proof}

\begin{problem}
	设 $f(x)$ 在 $[a,b]$ 上二阶连续可导，且 $f(a)=f(b)=0$，证明：
	\[
	\int_a^b f(x)\dd x=\frac12\int_a^b (x-a)(x-b)f''(x)\dd x,
	\]
	并且
	\[
	\left|\int_a^b f(x)\dd x\right|
	\leqslant \frac{(b-a)^3}{12}\max_{a\leqslant x\leqslant b}|f''(x)|.
	\]
\end{problem}

\begin{proof}
	由分部积分，
	\[
	\begin{aligned}
		\int_a^b (x-a)(x-b)f''(x)\dd x
		&=-\int_a^b (2x-a-b)f'(x)\dd x\\
		&=2\int_a^b f(x)\dd x.
	\end{aligned}
	\]
	故
	\[
	\int_a^b f(x)\dd x=\frac12\int_a^b (x-a)(x-b)f''(x)\dd x.
	\]
	令 $M=\max\limits_{a\leqslant x\leqslant b}|f''(x)|$，则
	\[
	\begin{aligned}
		\left|\int_a^b f(x)\dd x\right|
		&\leqslant \frac12 M\int_a^b (x-a)(b-x)\dd x\\
		&=\frac{M(b-a)^3}{12}.
	\end{aligned}
	\]
\end{proof}

\begin{problem}
	已知函数 $f(x)$ 在 $[a,b]$ 上三次可导，$f(a)=f'(a)=f(b)=0$，且 $|f'''(x)|\leqslant M$ 在 $[a,b]$ 上恒成立，证明：
	\[
	\left|\int_a^b f(x)\dd x\right|\leqslant \frac{M(b-a)^4}{72}.
	\]
\end{problem}

\begin{proof}
	取
	\[
	w(x)=(x-a)(x-b)^2.
	\]
	连续分部积分三次。因 $w(a)=w(b)=0$，$w'(b)=0$，且 $f(a)=f'(a)=f(b)=0$，有
	\[
	\begin{aligned}
		\int_a^b w(x)f'''(x)\dd x
		&=-\int_a^b w'(x)f''(x)\dd x\\
		&=\int_a^b w''(x)f'(x)\dd x\\
		&=-\int_a^b w'''(x)f(x)\dd x
		=-6\int_a^b f(x)\dd x.
	\end{aligned}
	\]
	于是
	\[
	\int_a^b f(x)\dd x=-\frac16\int_a^b (x-a)(x-b)^2 f'''(x)\dd x.
	\]
	从而
	\[
	\begin{aligned}
		\left|\int_a^b f(x)\dd x\right|
		&\leqslant \frac{M}{6}\int_a^b (x-a)(b-x)^2\dd x\\
		&=\frac{M}{6}\cdot\frac{(b-a)^4}{12}
		=\frac{M(b-a)^4}{72}.
	\end{aligned}
	\]
\end{proof}

\begin{problem}
	设 $f(x)$ 在 $[0,1]$ 上有连续的二阶导数，且
	\[
	f(0)=f(1)=f'(0)=0,\qquad f'(1)=1.
	\]
	证明：
	\[
	\int_0^1 \left(f''(x)\right)^2\dd x\geqslant4.
	\]
\end{problem}

\begin{proof}
	由条件，
	\[
	\int_0^1 f''(x)\dd x=f'(1)-f'(0)=1.
	\]
	又由分部积分，
	\[
	\int_0^1 x f''(x)\dd x
	=xf'(x)\bigg|_0^1-\int_0^1 f'(x)\dd x
	=1-f(1)+f(0)=1.
	\]
	于是
	\[
	\int_0^1 \left(x-\frac13\right)f''(x)\dd x=\frac23.
	\]
	由 Cauchy-Schwarz 不等式，
	\[
	\frac49
	\leqslant \int_0^1\left(x-\frac13\right)^2\dd x\int_0^1\left(f''(x)\right)^2\dd x
	=\frac19\int_0^1\left(f''(x)\right)^2\dd x.
	\]
	故
	\[
	\int_0^1\left(f''(x)\right)^2\dd x\geqslant4.
	\]
\end{proof}

\begin{problem}
	已知 $f(x)$ 为 $[0,2]$ 上的可微函数，且 $|f'(x)|\leqslant1$，$f(0)=f(2)=1$，证明：
	\[
	\int_0^2 f(x)\dd x>1.
	\]
\end{problem}

\begin{proof}
	当 $x\in[0,1]$ 时，由 $|f'|\leqslant1$ 得
	\[
	f(x)\geqslant f(0)-x=1-x.
	\]
	当 $x\in[1,2]$ 时，同理
	\[
	f(x)\geqslant f(2)-(2-x)=x-1.
	\]
	因此
	\[
	\int_0^2 f(x)\dd x
	\geqslant \int_0^1(1-x)\dd x+\int_1^2(x-1)\dd x=1.
	\]
	若等号成立，则必有
	\[
	f(x)=
	\begin{cases}
		1-x,&0\leqslant x\leqslant1,\\
		x-1,&1\leqslant x\leqslant2,
	\end{cases}
	\]
	这在 $x=1$ 处不可微，矛盾。故严格不等式成立。
\end{proof}

\begin{problem}
	已知 $f(x)\in C[a,b]$，且
	\[
	\int_a^b f(x)\dd x=\int_a^b x f(x)\dd x=0,
	\]
	则在 $(a,b)$ 内至少存在两个互异的点 $\xi_1,\xi_2$，使得
	\[
	f(\xi_1)=f(\xi_2)=0.
	\]
\end{problem}

\begin{proof}
	令
	\[
	F(x)=\int_a^x f(t)\dd t.
	\]
	由 $\displaystyle\int_a^b f(x)\dd x=0$，得 $F(a)=F(b)=0$。又
	\[
	0=\int_a^b x f(x)\dd x
	=xF(x)\bigg|_a^b-\int_a^b F(x)\dd x
	=-\int_a^b F(x)\dd x.
	\]
	令
	\[
	G(x)=\int_a^x F(t)\dd t.
	\]
	则 $G(a)=G(b)=0$。由 Rolle 定理，存在 $c\in(a,b)$，使得
	\[
	G'(c)=F(c)=0.
	\]
	于是 $F(a)=F(c)=0$，再由 Rolle 定理，存在 $\xi_1\in(a,c)$，使得
	\[
	F'(\xi_1)=f(\xi_1)=0.
	\]
	同理，由 $F(c)=F(b)=0$，存在 $\xi_2\in(c,b)$，使得
	\[
	F'(\xi_2)=f(\xi_2)=0.
	\]
	显然 $\xi_1\neq\xi_2$。
\end{proof}

\subsection*{Jensen 不等式}

\begin{proposition}
	（Jensen 不等式）设 $f\in\mathcal{R}[a,b]$，$g$ 为区间上的凸函数，且 $g(f(x))$ 可积，则
	\[
	g\left(\frac1{b-a}\int_a^b f(x)\dd x\right)
	\leqslant
	\frac1{b-a}\int_a^b g(f(x))\dd x.
	\]
\end{proposition}

\begin{proof}
	记
	\[
	m=\frac1{b-a}\int_a^b f(x)\dd x.
	\]
	由凸函数的切线性质，
	\[
	g(y)\geqslant g(m)+g'(m)(y-m).
	\]
	令 $y=f(x)$，得到
	\[
	g(f(x))\geqslant g(m)+g'(m)(f(x)-m).
	\]
	两边积分，得
	\[
	\int_a^b g(f(x))\dd x
	\geqslant (b-a)g(m)+g'(m)\left(\int_a^b f(x)\dd x-(b-a)m\right)
	=(b-a)g(m).
	\]
	整理即得 Jensen 不等式。
\end{proof}
